\section{Material Intervention and Model Readout}
\label{sec:method}

\paragraph{Changing the material representation.}
We construct two versions of each source scene. The original version retains
its MDL shader networks; the converted version uses UsdPreviewSurface networks.
The conversion traverses USD references, payloads, sublayers, variants, and
texture dependencies, preserving composition rather than flattening the scene.
Meshes and object transforms are retained. Matched camera settings expose the
resulting appearance difference, including changes in color, contrast, and
reflected light, to the model. Thus a material change can alter visual evidence
even when the spatial answer remains the same.

\paragraph{A spatial output has two components.}
A model first generates a response to an image and a query. A readout rule then
maps that response to a location. Let $I_i^O$ and $I_i^C$ denote paired images,
$q_i$ a query, and $B_i$ the target region. For a model $f$ and a prespecified
readout $r$, the decoded point is $p_i=r(f(I_i,q_i))$. A readout includes both
the response syntax and the coordinate frame. We measure the complete pipeline:
\begin{equation}
 H_i^v = \mathbf{1}\{p_i^v\text{ is a valid in-image point in } B_i\},
 \qquad
 \Delta_H = \frac{1}{N}\sum_{i=1}^{N}(H_i^C-H_i^O).
 \label{eq:hit}
\end{equation}
A missing or malformed point has $H_i^v=0$. Alongside $H$, we report
numeric-output coverage and the hit rate among
parsed numeric outputs. Coverage records whether a point can be decoded; the
end-to-end hit rate additionally requires it to lie inside the image and target.
All point displacements are
converted to image pixels before distances are computed.

\paragraph{Fixed-factor comparisons.}
The readout experiment crosses two requested frames (pixels and normalized
0--1000) with two syntaxes (JSON and two-line text), using the same image queries.
The six configurations are Gemma4-E4B in four-bit form
\citep{Google2026Gemma4}, Qwen2.5-VL-7B \citep{Bai2025Qwen},
Qwen3-VL-2B and 8B \citep{Bai2025Qwen3}, InternVL3-8B
\citep{Zhu2025InternVL}, and Molmo-7B-D \citep{Deitke2024Molmo}. The analysis preserves the
checkpoint, raw response, image digest, and prompt for every query. A common
0--1000 reinterpretation provides a second, uniformly applied diagnostic
alongside scoring in the requested frame.

\paragraph{Decision turnover.}
Mean differences can cancel. We therefore also count $O$-success/$C$-failure and
$O$-failure/$C$-success pairs, for both pointing and navigation. For multiple
views of one object, we compare the actual sets of successful views rather than
just their cardinalities. Uncertainty for point-hit differences is estimated by
resampling target clusters, keeping repeated views together. Experimental
details and paired render examples appear in Appendix~\ref{app:details};
Appendix~\ref{app:materials} compares the retained material cues.

\paragraph{From a point to a trajectory.}
For navigation, the material intervention is applied to the KuJiaLe scenes used
by the InternNav/DualVLN evaluation route. The language instruction, start state,
episode definition, and agent configuration are matched across conditions. We
record success rate (SR), success weighted by path length (SPL), navigation error
(NE), and trajectory length (TL), retaining episode IDs to connect aggregate
changes to individual outcomes.
